\begin{scanpage}{0164}{146}
In the case when the current filament passes through the point Q parallel to the half-plane's edge, the field at the point P is determined---in accordance with (20.11) and (20.12)---by the equations
\[
\left.\begin{aligned}E_z&=ik^2p_z[u(d,\varphi'-\varphi'')-u(d,\varphi'+\varphi'')]\sqrt{\frac{2\pi}{kR}}e^{i(kR-\pi/4)},\\
H_z&=ik^2m_z[u(d,\varphi'-\varphi'')+u(d,\varphi'+\varphi'')]\sqrt{\frac{2\pi}{kR}}e^{i(kR-\pi/4)}.
\end{aligned}\right\}\tag{24.15}
\]
With the absence of a half-plane, these sources create at the point P the field
\[
\left.\begin{aligned}E_z&=ik^2p_z\sqrt{\frac{2\pi}{kR}}e^{i(kR-\pi/4)}e^{-ikd\cos(\varphi'-\varphi'')},\\
H_z&=ik^2m_z\sqrt{\frac{2\pi}{kR}}e^{i(kR-\pi/4)}e^{-ikd\cos(\varphi'-\varphi'')}.
\end{aligned}\right\}\tag{24.16}
\]
Comparing Equations (24.15) and (24.16), we obtain the same Expressions (24.14) for the shading functions. Consequently, a spherical wave in the direction perpendicular to an ideally conducting half-plane is shaded by it the same as a cylindrical wave.

Let us note, however, that Expressions (24.14) are not equivalent to Expressions (21.22), (21.23) and (22.13), since the first represent the shading function by a half-plane of a wave from a single source, and the latter represent the shading function of an edge wave which we approximate by waves from two sources located on both sides of the corresponding half-plane. Since the shading functions of spherical and cylindrical waves are the same, the edge wave shading functions of a strip and a disk also will coincide.

Therefore, the approximation expressions for a field scattered by a disk which take account of secondary diffraction may be represented in the region $\varphi=-\pi/2$, $0<\vartheta<\pi/2$ (with $\zeta\gg1$) in the following form:
\end{scanpage}
